\chapter{Introduction}
\label{ch:introduction}

% ============================================================================
% This chapter introduces the research problem, motivation, and objectives
% ============================================================================

\section{Motivation}
\label{sec:motivation}

Sleep is a fundamental physiological process essential for human health and well-being. Sleep stage classification---the task of automatically identifying different stages of sleep from physiological signals---is critical for diagnosing sleep disorders, monitoring sleep quality, and advancing sleep research. Traditional manual sleep staging by trained experts is time-consuming, subjective, and expensive, making automated classification systems highly desirable.

Electroencephalography (EEG) provides rich information about brain activity during sleep and is the gold standard for sleep stage classification. However, the high dimensionality of EEG data, inter-subject variability, and class imbalance present significant challenges for machine learning models.

% TODO: Add more motivation - why is this important?
% - Clinical applications (sleep disorder diagnosis)
% - Research applications (large-scale sleep studies)
% - Cost and scalability benefits of automation

\section{Research Problem}
\label{sec:research_problem}

The primary research problem addressed in this thesis is:

\begin{quote}
\textit{How can we build an efficient, reproducible, and subject-independent sleep stage classification pipeline using machine learning models trained on EEG data?}
\end{quote}

This problem encompasses several key challenges:

\begin{enumerate}
    \item \textbf{Subject Independence:} Models must generalize to new subjects not seen during training, as evaluated by Leave-One-Subject-Out (LOSO) cross-validation.

    \item \textbf{Feature Engineering:} Extracting informative features from raw EEG signals that capture sleep stage characteristics across time and frequency domains.

    \item \textbf{Model Selection:} Comparing different machine learning models (XGBoost, Random Forest) for sleep stage classification.

    \item \textbf{Computational Efficiency:} The LOSO cross-validation with 128 subjects requires training 128 models per configuration, making computational efficiency critical.

    \item \textbf{Reproducibility:} Ensuring experiments are fully reproducible with proper cache invalidation when parameters change.
\end{enumerate}

\section{Research Objectives}
\label{sec:objectives}

The specific objectives of this thesis are:

\begin{enumerate}
    \item Implement a complete sleep stage classification pipeline from raw EEG data to predictions.

    \item Extract 149 time-domain and frequency-domain features from 6 EEG channels per 30-second epoch.

    \item Apply ANOVA-based feature selection to reduce dimensionality and improve model performance.

    \item Train and evaluate XGBoost and Random Forest models using LOSO cross-validation on 128 subjects from the BOAS dataset.

    \item Implement a two-tier caching system (feature cache + model cache) with SHA-256 fingerprint-based cache invalidation to accelerate iterative experimentation.

    \item Compare model performance across different feature selection configurations.

    \item Analyze the computational efficiency gains from the caching system.
\end{enumerate}

\section{Contributions}
\label{sec:contributions}

The main contributions of this thesis are:

\begin{enumerate}
    \item A complete, reproducible sleep stage classification pipeline implemented in Python with modular architecture.

    \item A two-tier caching system with fingerprint-based cache invalidation that reduces experiment time while ensuring reproducibility.

    \item Empirical comparison of XGBoost and Random Forest for subject-independent sleep stage classification on the BOAS dataset.

    \item Open-source implementation with comprehensive documentation for future research.
\end{enumerate}

\section{Thesis Structure}
\label{sec:structure}

The remainder of this thesis is organized as follows:

\begin{itemize}
    \item \textbf{Chapter 2} reviews related work in sleep stage classification, feature extraction, and machine learning models.

    \item \textbf{Chapter 3} describes the methodology, including the BOAS dataset, preprocessing pipeline, feature extraction, feature selection, and machine learning models.

    \item \textbf{Chapter 4} details the implementation of the pipeline, including the two-tier caching system and LOSO cross-validation framework.

    \item \textbf{Chapter 5} presents the experimental results, including model performance and caching efficiency.

    \item \textbf{Chapter 6} discusses the results, limitations, and implications.

    \item \textbf{Chapter 7} concludes the thesis and outlines directions for future work.
\end{itemize}
